\documentclass[twocolumn]{article}
\usepackage{amsmath,amssymb}
\begin{document}

\title{Horizontes de investigaci\'on:\\
  Tres preguntas abiertas del modelo FCC-DNLS}
\author{Serie Arm\'onica Collaboration}
\date{Junio 2026}
\maketitle

\begin{abstract}
  Todo marco te\'orico de frontera se valida por su poder de generar
  preguntas que antes no pod\'{\i}an formularse. El modelo FCC-DNLS
  resuelve problemas hist\'oricos (colapso cu\'antico sin par\'ametros
  libres, dimensionalidad del espacio, origen de la masa como
  auto-atrapamiento), pero tambi\'en abre tres horizontes
  radicalmente nuevos. Formalizamos cada uno con predicciones
  comprobables y programas de investigaci\'on.
\end{abstract}

\section{Introducci\'on}

El modelo unificado FCC-DNLS (v3) ha establecido:
\begin{itemize}
  \item Punto fijo espectral en 3D: ancho $2.96$, energ\'{\i}a $-0.056$
  \item Purificaci\'on $S(t)\to0$, rango $\rho: 68\to1$
  \item Umbral de auto-atrapamiento en Norma $=4$
  \item $\beta = (1-\eta) + \varepsilon_0/2 = 0.30119$, derivado de la
    fracci\'on de empaquetamiento FCC $\eta = \pi/(3\sqrt2)$
  \item Kernel de colapso $K(\Delta n)$ con soporte compacto
  \item Ecuaci\'on funcional de Epstein verificada exactamente
\end{itemize}

Sobre esta base, el modelo abre tres preguntas que exploramos a
continuaci\'on.

\section{Pregunta A: \textquestiondown Es el universo temprano un
  efecto de auto-interferencia espectral?}

\subsection{El problema}

El telescopio James Webb (JWST) detecta galaxias masivas
($M_\star \sim 10^{10-11}M_\odot$) a redshifts $z>10$, cuando el
universo ten\'{\i}a menos de 500~millones de a\~nos. En
$\Lambda$CDM se requieren eficiencias de formaci\'on estelar
$\epsilon > 1$.

\subsection{Respuesta del modelo}

El universo observable ser\'{\i}a un solit\'on espectral en una red
FCC infinita. El corrimiento al rojo no es la expansi\'on FLRW sino
el gradiente espectral del perfil del solit\'on:

\begin{equation}
  \frac{\delta\lambda}{\lambda} \approx \frac{r^2}{2\xi^2},
\end{equation}

donde $r$ es la distancia al centro y $\xi \approx 2.96$ el ancho
caracter\'{\i}stico. Un objeto al borde del solit\'on ($r\sim\xi$)
sufre un corrimiento $\delta\lambda/\lambda\sim0.5$, interpretado
en FLRW como $z\sim1$.

\subsection{Predicciones falsables}

\begin{enumerate}
  \item[A1.] La relaci\'on distancia-redshift se desv\'{\i}a de FLRW
    para $z>6$ siguiendo una ley cuadr\'atica. Falsable con $H(z)$
    de SNIa + JWST.
  \item[A2.] Las galaxias $z>10$ tienen morfolog\'{\i}as
    ``espectralmente replicadas'' de galaxias $z\sim2$.
  \item[A3.] El CMB tiene una modulaci\'on dipolar residual del
    perfil del solit\'on, no atribuible al movimiento solar.
\end{enumerate}

\section{Pregunta B: \textquestiondown Es la no-localidad cu\'antica
  una ilusi\'on de la base espacial?}

\subsection{El problema}

Las violaciones de Bell (Aspect 1982, Hensen 2015) demuestran
correlaciones cu\'anticas que no admiten variables ocultas locales.
Se interpretan como ``acci\'on fantasmal a distancia''.

\subsection{Respuesta del modelo}

El colapso es local en la base espectral (autovectores del
Laplaciano). La transformaci\'on entre la base espectral y la base
real es la transformada de Fourier. Un operador multiplicativo en
$k$-espacio es una convoluci\'on en $x$-espacio:

\begin{equation}
  K(\Delta n) = \sum_k \sigma(\lambda_k)
    \frac{\varepsilon_0}{\varepsilon_0+|c_k|^2}
    e^{i k\cdot\Delta n}.
\end{equation}

El n\'ucleo $K(\Delta n)$ tiene soporte compacto: $K\approx0$ para
$|\Delta n| > 1$ sitio (verificado: $1/e$ a 1.0~sitios). La
``no-localidad'' cu\'antica emerge porque las mediciones se
realizan en base real, pero el colapso es local en base espectral.
La transformada de Fourier proyecta esta localidad espectral como
correlaci\'on instant\'anea en espacio real.

\subsection{Predicciones falsables}

\begin{enumerate}
  \item[B1.] La violaci\'on de CHSH decae con la distancia como
    $K(\Delta n)$. Falsable con experimentos Bell en redes \'opticas
    FCC con separaci\'on variable.
  \item[B2.] Existe una descripci\'on local oculta en $k$-espacio:
    $\lambda_k = \{\arg(c_k), |c_k|^2\}$ reproduce las estad\'{\i}sticas
    cu\'anticas.
  \item[B3.] El colapso espectral induce una fase geom\'etrica
    medible en interferometr\'{\i}a, proporcional a $\beta\varepsilon_0$.
\end{enumerate}

\section{Pregunta C: \textquestiondown Cu\'al es el papel de los
  n\'umeros primos en el espacio-tiempo discreto?}

\subsection{El problema}

La funci\'on zeta de Riemann $\zeta(s)$ codifica la distribuci\'on
de n\'umeros primos. La zeta de Epstein $\zeta_{D_3}(s)$ para la
red $D_3$ (la red FCC) satisface $\xi(3/2-s)=\xi(s)$.

\subsection{Respuesta del modelo}

El Laplaciano FCC genera autovalores $\lambda_k$ cuya distribuci\'on
determina: (1) la dispersi\'on de ondas, (2) la funci\'on zeta
espectral $\zeta_\Delta(s)$, (3) la densidad de estados
$\rho(\lambda)$, (4) el umbral de auto-atrapamiento, (5) la
estabilidad del atractor.

La funci\'on zeta espectral y la zeta de Epstein est\'an conectadas
por el teorema de Friedli-Karlsson:

\begin{equation}
  \xi(s) = \pi^{-s}\Gamma(s)\,\zeta_{D_3}(s),\quad
  \xi(3/2-s)=\xi(s).
\end{equation}

El error es $0.000000$ para $N=4,6,8,10$: la conexi\'on es exacta,
no asint\'otica. Los autovalores del Laplaciano FCC para $k$
peque\~nos se aproximan a $|k|^2$, donde $|k|^2$ recorre las
representaciones de enteros como suma de tres cuadrados (coeficientes
de la funci\'on theta de Jacobi $\vartheta_3(q)^3$).

\subsection{Predicciones falsables}

\begin{enumerate}
  \item[C1.] El error en $\xi(3/2-s)=\xi(s)$ escala como
    $O(N^{-\alpha})$ con $\alpha \approx 1/2$.
  \item[C2.] El umbral de auto-atrapamiento (Norma $=4$) corresponde
    a un autovalor aritm\'eticamente distinguido del Laplaciano FCC.
  \item[C3.] Los 33 canales de baja frecuencia se organizan en
    multipletes del grupo de Weyl de $D_3$, con degeneraciones dadas
    por coeficientes de representaciones enteras.
\end{enumerate}

\section{Conclusiones}

Cada programa genera predicciones falsables expl\'{\i}citas con
tecnolog\'{\i}a actual (redes \'opticas FCC, JWST, interferometr\'{\i}a
at\'omica). El modelo FCC-DNLS no es una especulaci\'on
filos\'ofica: es un marco computacional y matem\'atico que puede
ser verificado o refutado en un horizonte de 3--5 a\~nos.

\begin{thebibliography}{9}
\bibitem{friedli2017} F.~Friedli and A.~Karlsson, Tohoku Math. J.
  {\bf 69}(4), 585 (2017).
\bibitem{tissot2024} B.~Tissot, H.~Ribeiro, and F.~Marquardt,
  Phys. Rev. Research {\bf 6}, 023015 (2024).
\bibitem{westhoff2025} P.~Westhoff, S.~Paeckel, and M.~Moroder,
  Phys. Rev. A {\bf 112}, L061304 (2025).
\bibitem{epstein1903} P.~Epstein, Math. Ann. {\bf 56}, 615 (1903).
\bibitem{aspect1982} A.~Aspect, J.~Dalibard, and G.~Roger,
  Phys. Rev. Lett. {\bf 49}, 1804 (1982).
\bibitem{hensen2015} B.~Hensen {\it et al.}, Nature {\bf 526},
  682 (2015).
\end{thebibliography}

\end{document}
